\chapterwithopening{Future Work}{chap:future_work}{
  \chaptercontentsitem{sec:future_work_introduction}
  \chaptercontentsitem{sec:scalability_improvements}
  \chaptercontentsitem{sec:user_feedback_usability_refinement}
  \chaptercontentsitem{sec:mobile_companion_expansion}
  \chaptercontentsitem{sec:testing_quality_assurance_expansion}
  \chaptercontentsitem{sec:localization_accessibility}
  \chaptercontentsitem{sec:advanced_ide_assessment_enhancements}
  \chaptercontentsitem{sec:future_work_summary}
}

\section{Introduction}\label{sec:future_work_introduction}

EduVerse already demonstrates a broad and coherent educational workspace, but
the current implementation also reveals several realistic directions for future
improvement. These directions do not represent missing fundamentals of the
implemented system. Rather, they reflect the next stage of refinement expected
when a multi-organization academic platform moves from successful feature
integration toward wider adoption, heavier usage, and more formal evaluation.

The most immediate future work concerns scalability and structured user
feedback. Additional improvements remain valuable as well, especially in the
areas of mobile expansion, testing maturity, localization, accessibility, and
advanced educational tooling.

\section{Scalability Improvements}\label{sec:scalability_improvements}

The first major future-work direction is scalability. EduVerse already manages
organizations, classes, files, results, notifications, live sessions, and
several feature-gated tools, so future growth will depend on how efficiently
these services continue to operate as user counts, class counts, and stored
content increase.

In practical terms, scalability work should focus on heavier roster loading,
message volume, results aggregation, large material libraries, and concurrent
live-session usage. The current implementation already separates large binary
files from relational records, which is a strong foundation. Future work can
extend that foundation through deeper performance monitoring, query
optimization, caching strategy review, and operational tuning based on real
usage patterns.

\section{User Feedback and Usability Refinement}
\label{sec:user_feedback_usability_refinement}

The second major future-work direction is systematic feedback collection from
real users. EduVerse serves Organization Administrators, Teachers, and
Students, and each of these groups experiences the platform differently. As the
system is used more widely, direct feedback from those groups should guide
refinement of navigation clarity, role switching, classroom workflow
understandability, notification usefulness, and help content quality.

This kind of refinement is especially important in a platform that combines
many related academic tools within one workspace. The technical integration is
already a strength, but long-term usability depends on whether the integrated
experience remains clear and low-friction for different users under everyday
academic conditions.

\section{Mobile Companion Expansion}\label{sec:mobile_companion_expansion}

The project documentation references a separate mobile companion repository, but
that codebase is not part of the current implementation evidence reviewed in
this thesis. Future work should therefore focus on expanding that companion in
a controlled way while keeping the web application as the primary source of
full administrative depth.

The most important objective is not simple screen duplication. It is
contract-level consistency between the mobile client and the shared API layer so
that authentication, role resolution, class access, notifications, materials,
assignments, and live-session entry remain aligned across clients. Future work
in this area should include stronger mobile/web contract validation and careful
selection of which workflows genuinely benefit from mobile-first interaction.

\section{Testing and Quality Assurance Expansion}
\label{sec:testing_quality_assurance_expansion}

The current repository shows meaningful logic-level test coverage, but the next
stage of quality assurance should add broader end-to-end workflow testing.
Cross-role scenarios such as invite acceptance, class entry, submission,
grading, result release, and public join flows would benefit from automated
browser-level verification.

This work is particularly important because EduVerse combines access control,
feature gating, and multi-step academic workflows. As the platform grows,
end-to-end testing will help reduce regression risk in ways that unit and
helper-level tests alone cannot fully cover. The same expansion should also
include stronger integration checks between the main web client and the mobile
companion direction.

\section{Localization and Accessibility}
\label{sec:localization_accessibility}

The implementation evidence shows that language switching is not yet complete
for end users. Future work should therefore address localization if the system
is expected to support multilingual academic environments. This includes both
translation of interface text and review of how role-specific terminology,
notifications, and help content are presented across languages.

Accessibility should advance in parallel with localization. A platform intended
for long-term educational use should be reviewed for keyboard access, visual
clarity, screen-reader compatibility, focus management, and the accessibility
of assessments, live-session controls, and coding interfaces.

\section{Advanced IDE and Assessment Enhancements}
\label{sec:advanced_ide_assessment_enhancements}

The current browser-based IDE is appropriate for guided educational coding
practice, but future work could extend it through a secure backend execution
sandbox if real command execution or compiled-language support becomes a
requirement. Such an enhancement would need strong isolation and careful
resource control so that flexibility does not compromise security.

Assessment workflows can also be extended further if institutional needs become
more demanding. The current implementation already supports event-based exam
integrity monitoring, passcodes, retakes, and explicit result release. Future
work may explore stronger integrity controls or richer assessment analytics, but
such additions should remain proportionate to actual educational requirements
rather than being treated as mandatory by default.

\section{Chapter Summary}\label{sec:future_work_summary}

Future work for EduVerse should focus first on scalability improvements and
structured feedback from real users, because those two areas most directly
shape whether the integrated platform can mature beyond its current successful
implementation scope. Additional priorities include mobile companion
expansion, end-to-end testing, localization, accessibility, and carefully
governed enhancement of the IDE and assessment subsystems.
